\subsection{实施条件与阶段稳定性}
\subsubsection{成本与执行时点}
表\ref{tab:implementation_core}保留基准费率、较高费率和延迟执行三个代表情景。成本变化沿用同一目标权重路径；延迟实验将原决策顺延一个交易日执行，并重新核算漂移持仓、交易量与净收益。
\begin{table}[htbp]\centering
\caption{成本与执行条件的核心比较}\label{tab:implementation_core}
\begin{tabular}{lr}\toprule
情景 & 净年化收益\\\midrule
基准：单边\txCostBps{} bps & \globalNetReturn\\
单边成本20 bps & \diagCostTwentyReturn\\
延后一交易日执行 & \diagDelayReturn\\\bottomrule
\end{tabular}
\end{table}
\FloatBarrier
较高费用降低了净收益；延后一日执行时，年化波动为\diagDelayVol，最大回撤为\diagDelayDD，与基准结果接近。这些检验说明结果在所考察的实施变化下未出现明显反转，同时表明交易成本仍是配置效果的一部分。完整费率、风险指标和均值半衰期敏感性见表\ref{tab:diagnostic_robustness}。

\subsubsection{不同时间阶段}
表\ref{tab:diagnostic_subperiods}使用预先固定的三个连续子区间，比较主模型与两条风险匹配辅助路径。区间划分覆盖不同市场阶段，2026年只计入截至样本终点的观察。每个子区间重新以首日净值为基期计算绩效，回撤使用区间内历史高点。

\begin{table}[htbp]\centering
\caption{固定子区间中的收益、夏普与回撤}\label{tab:diagnostic_subperiods}
\resizebox{\textwidth}{!}{\begin{tabular}{llrrr}
\toprule
子区间 & 配置 & 净年化收益 & 夏普 & 最大回撤\\\midrule
2018--2020 & Global RRP & 5.95\% & 1.506 & -6.80\%\\
2018--2020 & Equal Weight（风险匹配） & 3.46\% & 1.183 & -4.53\%\\
2018--2020 & 60/40 Benchmark（风险匹配） & 3.89\% & 1.265 & -4.29\%\\
2021--2023 & Global RRP & 3.82\% & 1.204 & -2.83\%\\
2021--2023 & Equal Weight（风险匹配） & 2.20\% & 0.955 & -2.19\%\\
2021--2023 & 60/40 Benchmark（风险匹配） & 1.34\% & 0.569 & -3.22\%\\
2024--2026 & Global RRP & 8.34\% & 2.199 & -2.46\%\\
2024--2026 & Equal Weight（风险匹配） & 3.98\% & 1.474 & -1.98\%\\
2024--2026 & 60/40 Benchmark（风险匹配） & 3.54\% & 1.266 & -2.32\%\\
\bottomrule
\end{tabular}}
\end{table}

主模型在三个子区间的净年化收益均高于两条风险匹配辅助路径，2021--2023年的收益水平较低，2024--2026年较高。回撤关系则随区间变化：在2018--2020年，主模型的区间最大回撤更大；在2021--2023年，其最大回撤较小。由此可见，平均收益优势与阶段损失控制需要分别评价。子区间证据刻画了结果的时间分布，宏观增长、利率与通胀风险可作为理解资产暴露的经济背景，具体因果归因仍以相应因子数据为条件。
\FloatBarrier


以上证据表明，两项惩罚在当前样本中形成了可解释的收益风险调整，效果仍随时间阶段和实施条件变化。配对区块抽样及其区间见附录\ref{app:empirical}，正文以核心对照和实际取舍评价模型有效性。
